% =============================================================================
% Pregunta 3 — Aplicación de DDD y Bounded Context en un dominio complejo
% Autor: Ernesto A. Soler Choong (Team Bales) · Issue #5
%
% Fragmento para incluir con \input{} en el documento final. Requiere en el
% preámbulo: babel (spanish, es-nolists), csquotes, biblatex (style=apa), booktabs,
% tabularx, enumitem, listings (estilo «codigo»), xcolor (azul, naranja, gris)
% y tikz (positioning, arrows.meta, calc). Ver p3-standalone.tex.
% =============================================================================

\section{Aplicación de DDD y \textit{Bounded Context} en un dominio complejo}
\label{sec:p3}

\subsection{Planteamiento}
\label{sec:p3-planteamiento}

Una empresa de seguros necesita un sistema para gestionar pólizas de automóviles que
interactúe con el departamento de siniestros, con el sistema de facturación de primas y
con la aplicación que usan los ajustadores para evaluar daños en el lugar del accidente.
Cada área usa su propia terminología: para un actuario la póliza es un contrato con
cláusulas; para facturación, una fuente de ingresos; para un ajustador, un número de
referencia. La pregunta pide explicar:

\begin{enumerate}[label=(\alph*)]
  \item cómo dividir el sistema en contextos delimitados (\textit{bounded contexts});
  \item cómo evitar que el modelo de \enquote{Póliza} se contamine con detalles de
        siniestros o de facturación;
  \item cómo se comunicarían esos contextos de manera desacoplada, y
  \item cómo se relaciona la solución con una arquitectura orientada al dominio (Onion o
        Clean) y con el \enquote{Adaptador de Servicio} de la Guía de Arquitectura N-Capas.
\end{enumerate}

La solución ingenua sería construir un único modelo de \enquote{Póliza} compartido por
todas las áreas. Ese modelo acumularía cláusulas, planes de cobro y datos de peritaje, y
cualquier cambio pedido por un área afectaría a todas las demás. Domain-Driven Design
(DDD) propone lo contrario: expresar el modelo en el lenguaje de los expertos del
negocio, el llamado \emph{Lenguaje Ubicuo}, y está indicado precisamente para
aplicaciones empresariales complejas \parencite[21]{delatorre2010}. Desde esta
perspectiva, que \enquote{póliza} signifique tres cosas distintas no es un defecto que
haya que corregir, sino la señal de que existen varios modelos legítimos.

\subsection{División del sistema en contextos delimitados}
\label{sec:p3-division}

El contexto delimitado es un concepto propuesto por Evans (como se citó en
\cite[cap.~7]{richardsford2020}): todo lo relativo a un dominio es visible dentro del
contexto y opaco fuera de él. En vez de crear una única clase compartida por toda la
organización, cada dominio define la suya y las diferencias se concilian en los puntos de
integración \parencite[cap.~7]{richardsford2020}. Para trazar las fronteras en el caso de
la aseguradora se aplican cuatro criterios:

\begin{enumerate}[label=\textbf{C\arabic*.}, leftmargin=2.2em]
  \item \textbf{Lenguaje.} Cada área cuyos expertos usan un vocabulario propio es
        candidata a un contexto propio. El criterio se deriva de que el modelo debe
        expresarse en el Lenguaje Ubicuo de los expertos \parencite[21]{delatorre2010} y
        de que cada entidad funciona mejor dentro de un contexto localizado
        \parencite[cap.~7]{richardsford2020}.
  \item \textbf{Partición por dominio, no técnica.} Si los componentes se organizan por
        capas técnicas, cada flujo de negocio queda repartido entre todas ellas. Si se
        organizan por dominio, la estructura refleja mejor los cambios que realmente
        ocurren \parencite[cap.~8]{richardsford2020}.
  \item \textbf{Granularidad.} Cada contexto debe tener un propósito funcionalmente
        cohesivo. Las entidades que necesitan cooperar en una misma transacción sugieren
        una frontera común. Si dos contextos necesitan comunicarse en exceso, puede
        considerarse fusionarlos \parencite[cap.~17]{richardsford2020}.
  \item \textbf{Características operativas.} Un \emph{cuanto arquitectónico} es un
        artefacto desplegable de forma independiente, con alta cohesión funcional y
        connascencia síncrona. Una llamada síncrona obliga a que las dos partes tengan
        las mismas características operativas mientras dura la llamada
        \parencite[cap.~7]{richardsford2020}. Por eso conviene no acoplar de forma
        síncrona partes con exigencias operativas muy distintas.
\end{enumerate}

Con estos criterios se proponen cuatro contextos delimitados, resumidos en la
Tabla~\ref{tab:p3-contextos}. Cada uno tiene su propio significado de \enquote{Póliza}.

\begin{table}[htbp]
  \centering
  \caption{Contextos delimitados propuestos y significado de \enquote{Póliza} en cada uno.}
  \label{tab:p3-contextos}
  \small
  \renewcommand{\arraystretch}{1.25}
  \begin{tabularx}{\textwidth}{@{}>{\raggedright\arraybackslash}p{2.6cm}
                                  >{\raggedright\arraybackslash}X
                                  >{\raggedright\arraybackslash}X
                                  >{\raggedright\arraybackslash}X@{}}
    \toprule
    \textbf{Contexto} & \textbf{Responsabilidad} & \textbf{Lenguaje propio}
      & \textbf{Qué es \enquote{Póliza}} \\
    \midrule
    Suscripción y Pólizas (actuarios)
      & Diseñar, emitir, modificar, renovar, suspender y cancelar pólizas
      & cobertura, cláusula, riesgo, prima técnica, vigencia
      & Contrato con coberturas y cláusulas; raíz del agregado principal \\
    Facturación de Primas
      & Emitir recibos, cobrar primas y gestionar impagos
      & recibo, plan de pagos, vencimiento, impago
      & Fuente de ingresos: cuenta de cobro asociada a un número de póliza \\
    Siniestros
      & Abrir y tramitar expedientes y decidir indemnizaciones
      & expediente, declaración, reserva, indemnización
      & Cobertura aplicable en la fecha del siniestro \\
    Peritaje (ajustadores)
      & Recibir solicitudes de peritaje, asignar ajustadores, evaluar daños en el lugar
        del accidente y emitir el informe pericial
      & inspección, daño, valoración, informe pericial
      & Número de póliza usado como mera referencia del vehículo asegurado \\
    \bottomrule
  \end{tabularx}
\end{table}

\paragraph{Decisión de diseño: Peritaje como contexto propio.}
La decisión menos evidente es separar Peritaje de Siniestros, aunque ambos participan en el
mismo proceso. Hay tres argumentos a favor:
\begin{itemize}
  \item Su lenguaje es distinto: inspección y valoración frente a expediente y reserva
        (criterio C1).
  \item No comparten transacción: el informe pericial es un producto terminado que se
        entrega a Siniestros (criterio C3).
  \item Sus características operativas difieren: Peritaje es una aplicación móvil usada
        en campo, posiblemente con conectividad intermitente, y acoplarla de forma
        síncrona a Siniestros le impondría las mismas exigencias (criterio C4).
\end{itemize}
La alternativa de integrarlo en Siniestros sería razonable en una organización pequeña.
Ningún estilo de partición es intrínsecamente más correcto que otro
\parencite[cap.~8]{richardsford2020}, y la granularidad adecuada se alcanza iterando
\parencite[cap.~17]{richardsford2020}.

Conviene distinguir esta división de los \emph{Módulos} que propone la Guía de
Arquitectura N-Capas: particiones verticales del modelo de dominio por áreas de negocio,
con alta cohesión y comunicadas mediante interfaces \parencite[71--72]{delatorre2010}. Un
contexto delimitado va un paso más allá. No es un módulo de un modelo compartido: es un
modelo propio, con sus propios datos.

\subsection{Cómo evitar la contaminación del modelo de Póliza}
\label{sec:p3-contaminacion}

El aislamiento del modelo se consigue combinando cinco mecanismos.

\begin{enumerate}[label=\textbf{M\arabic*.}, leftmargin=2.2em]
  \item \textbf{Un modelo de Póliza por contexto.} No existe una clase \enquote{Póliza}
        compartida. Cada contexto incluye sus propias clases y esquemas, y nunca se acopla
        a la definición de clases ni a la base de datos de otro contexto
        \parencite[cap.~17]{richardsford2020}. Solo Suscripción conoce cláusulas y
        coberturas. Siniestros maneja una \emph{cobertura aplicable} y Facturación una
        \emph{cuenta de cobro}.

  \item \textbf{Aceptar la duplicación accidental.} \textcite[154--155]{martin2017}
        distingue la duplicación verdadera, en la que todo cambio en una copia exige el
        mismo cambio en las demás, de la accidental, en la que las copias evolucionan por
        caminos distintos. Unificar la duplicación accidental crea un acoplamiento difícil
        de deshacer. Ambos contextos, Suscripción y Siniestros, manejan datos de
        cobertura (límites, franquicia, vigencia). Suscripción los modela como cláusulas
        de un contrato que se negocian y renuevan; Siniestros solo necesita una
        instantánea para resolver un expediente. Las dos estructuras cambiarán por
        razones distintas, así que su parecido es accidental. Cuando el objetivo es el
        desacoplamiento, se prefiere la duplicación a la reutilización
        \parencite[cap.~17]{richardsford2020}. El número de póliza, en cambio, no es una
        duplicación: es un identificador compartido que sirve de referencia entre
        contextos (mecanismo M3).

  \item \textbf{Fronteras de agregado y referencias por identidad.} Un agregado es un
        conjunto de objetos asociados que se trata como una unidad frente a los cambios
        de datos. Tiene una entidad raíz como único punto de entrada, y desde fuera solo
        puede referenciarse esa raíz, lo que asegura la integridad de los datos
        \parencite[172]{delatorre2010}. En Suscripción, la Póliza es la raíz de un
        agregado que contiene sus coberturas y cláusulas. Por analogía, la propuesta
        aplica la misma regla entre contextos: los demás contextos se refieren a la póliza
        solo por su número de póliza, nunca por su estructura interna.

  \item \textbf{Datos aislados.} Cada contexto tiene su propio almacenamiento y no se
        integra a través de esquemas compartidos \parencite[cap.~17]{richardsford2020}.

  \item \textbf{Traducción en la frontera.} La información que llega de otro contexto
        entra solo a través de un componente que la convierte al modelo local. Los
        Agentes de Servicios externos de la Guía cumplen este papel: entre otras tareas,
        realizan mapeos básicos entre el formato de datos del servicio externo y el que usa
        la propia aplicación \parencite[108, 155]{delatorre2010}. Un adaptador cumple la
        misma función en la Arquitectura Hexagonal: la aplicación interactúa con él sin
        conocer la naturaleza de lo que hay al otro lado
        \parencite{cockburn2005}.\footnote{En la terminología original de DDD, este
        mecanismo de traducción defensiva se conoce como \textit{Anti-Corruption Layer}.
        El término se menciona solo como orientación, porque no procede de las fuentes
        consultadas.}
\end{enumerate}

\subsection{Comunicación desacoplada entre contextos}
\label{sec:p3-comunicacion}

\textcite[cap.~14]{richardsford2020} distinguen dos modelos de interacción. El modelo
basado en peticiones es el adecuado para consultas de datos estructuradas que necesitan
certeza y control del flujo. El modelo basado en eventos es el adecuado para acciones que
requieren capacidad de respuesta, escalabilidad y adaptabilidad, a cambio de admitir solo
consistencia eventual. La propuesta combina ambos (Figura~\ref{fig:p3-mapa}).

\begin{figure}[htbp]
  \centering
  \begin{tikzpicture}[
      ctx/.style={draw=azul, fill=azul!6, rounded corners=4pt, thick,
                  minimum width=5.0cm, minimum height=2.2cm, align=center, font=\small},
      ev/.style={-{Stealth[length=2.4mm]}, thick, black!75},
      sync/.style={{Stealth[length=2.4mm]}-{Stealth[length=2.4mm]}, thick, dashed, naranja},
      lbl/.style={font=\scriptsize, align=center, fill=white, inner sep=1.5pt}
    ]
    \node[ctx] (sus) at (0,0)     {\textbf{Suscripción y Pólizas}\\[2pt]
                                   Póliza = contrato con\\ coberturas y cláusulas};
    \node[ctx] (fac) at (9.6,0)   {\textbf{Facturación de Primas}\\[2pt]
                                   Póliza = cuenta de cobro};
    \node[ctx] (sin) at (0,-5.4)  {\textbf{Siniestros}\\[2pt]
                                   Póliza = cobertura aplicable\\ en la fecha del siniestro};
    \node[ctx] (per) at (9.6,-5.4) {\textbf{Peritaje}\\[2pt]
                                   Póliza = número de póliza,\\ usado como mera referencia};

    \draw[ev] ([yshift=5mm]sus.east) -- node[lbl,above=1pt]
      {PólizaEmitida, PólizaRenovada,\\ PólizaModificada, PólizaSuspendida,\\ PólizaCancelada}
      ([yshift=5mm]fac.west);
    \draw[ev] ([yshift=-5mm]fac.west) -- node[lbl,below=1pt]{PagoVencido}
      ([yshift=-5mm]sus.east);
    \draw[ev] ([yshift=5mm]sin.east) -- node[lbl,above=1pt]{PeritajeSolicitado}
      ([yshift=5mm]per.west);
    \draw[ev] ([yshift=-5mm]per.west) -- node[lbl,below=1pt]{PeritajeCompletado}
      ([yshift=-5mm]sin.east);
    \draw[ev] ([xshift=-12mm]sin.north) -- node[lbl,left=1pt]{SiniestroLiquidado}
      ([xshift=-12mm]sus.south);
    \draw[sync] ([xshift=12mm]sin.north) -- node[lbl,right=1pt]
      {ConsultarCobertura\\(consulta síncrona\\a la API publicada)} ([xshift=12mm]sus.south);

    \draw[ev]   (-2.4,-7.3) -- ++(1.1,0) node[right,font=\scriptsize]
      {evento de dominio asíncrono (bróker, publicar/suscribir)};
    \draw[sync] (7.4,-7.3) -- ++(1.1,0) node[right,font=\scriptsize]
      {consulta síncrona};
  \end{tikzpicture}
  \caption{Mapa de contextos propuesto. Los eventos se publican en un bróker de
           mensajes; la flecha discontinua es la única interacción síncrona.}
  \label{fig:p3-mapa}
\end{figure}

\paragraph{Eventos de dominio asíncronos como mecanismo por defecto.}
Cada contexto anuncia lo que ha hecho y los demás reaccionan, sin un coordinador central.
Es la topología \textit{broker}, en la que los procesadores publican eventos en canales
de un intermediario ligero mediante publicación y suscripción
\parencite[cap.~14]{richardsford2020}. La coreografía, es decir, la ausencia de un
mediador global, respeta la filosofía del contexto delimitado
\parencite[cap.~17]{richardsford2020}. La difusión de eventos es probablemente el mayor
grado de desacoplamiento posible, porque el productor desconoce quién recibe el mensaje y
qué hace con él \parencite[cap.~14]{richardsford2020}. La Tabla~\ref{tab:p3-eventos}
detalla los flujos principales. No pretende ser exhaustiva: otros casos, como la
regularización de un pago o la cancelación de un peritaje, seguirían el mismo patrón.

\begin{table}[htbp]
  \centering
  \caption{Interacciones entre contextos.}
  \label{tab:p3-eventos}
  \small
  \renewcommand{\arraystretch}{1.2}
  \begin{tabularx}{\textwidth}{@{}>{\raggedright\arraybackslash}p{3.6cm}
                                  >{\raggedright\arraybackslash}p{2.3cm}
                                  >{\raggedright\arraybackslash}p{2.3cm}
                                  >{\raggedright\arraybackslash}X@{}}
    \toprule
    \textbf{Mensaje} & \textbf{Publica} & \textbf{Consume} & \textbf{Reacción del consumidor} \\
    \midrule
    PólizaEmitida, PólizaRenovada, PólizaModificada, PólizaSuspendida, PólizaCancelada
      & Suscripción & Facturación
      & Crea, ajusta, suspende o cierra el plan de cobro; el evento incluye el número de
        póliza, la prima y el calendario de pagos \\
    PagoVencido & Facturación & Suscripción
      & Decide, según las cláusulas del contrato, si suspende la póliza \\
    PeritajeSolicitado & Siniestros & Peritaje
      & Asigna un ajustador; el evento incluye la referencia de póliza y los datos del
        vehículo y del lugar recogidos en la declaración del siniestro \\
    PeritajeCompletado & Peritaje & Siniestros & Incorpora la valoración de daños al expediente \\
    SiniestroLiquidado & Siniestros & Suscripción
      & Actualiza el historial de siniestralidad que se usa en las renovaciones \\
    \addlinespace
    ConsultarCobertura (síncrona) & Siniestros pregunta & Suscripción responde
      & Verifica la cobertura vigente en la fecha del siniestro al abrir el expediente \\
    \bottomrule
  \end{tabularx}
\end{table}

Dos decisiones de la tabla merecen justificación:
\begin{itemize}
  \item \textbf{Cada contexto conserva la propiedad de sus decisiones.} Facturación no
        cambia el estado de una póliza: publica \texttt{PagoVencido} y es Suscripción
        quien decide la suspensión aplicando las cláusulas del contrato. Así, las reglas
        contractuales no se filtran a Facturación y los datos de cobro no se filtran a la
        Póliza.
  \item \textbf{Peritaje recibe lo que necesita.} El evento \texttt{PeritajeSolicitado}
        incluye la referencia de póliza y los datos del vehículo y del lugar que el
        asegurado aportó al declarar el siniestro. Así, el ajustador no depende de una
        consulta síncrona desde el lugar del accidente.
\end{itemize}

\paragraph{Extensibilidad.}
Que cada contexto anuncie lo que hizo, aunque hoy nadie lo escuche, facilita la
extensibilidad \parencite[cap.~14]{richardsford2020}. Por ejemplo, un futuro contexto de
detección de fraude podría suscribirse a los eventos de Siniestros sin modificar ese
contexto.

\paragraph{Consulta síncrona solo donde se necesita una respuesta inmediata.}
Al abrir un expediente, Siniestros necesita saber, con la información más actualizada
posible, si existía cobertura en la fecha del accidente. Es una consulta de datos determinista, el caso para el que se
recomienda el modelo basado en peticiones \parencite[cap.~14]{richardsford2020}. En el
diseño se implementa como una consulta síncrona a la API publicada de Suscripción
(Listado~\ref{lst:p3-agente}). Dentro de una arquitectura dirigida por eventos, la
alternativa equivalente sería la mensajería petición--respuesta
\parencite[cap.~14]{richardsford2020}. Se prefiere la API porque, para una única
interacción, evita añadir colas de petición y de respuesta. Otra alternativa sería
mantener en Siniestros una réplica local de las coberturas, alimentada por los eventos de
póliza. Eliminaría la dependencia temporal, pero la verificación también pasaría a ser
eventualmente consistente; el diseño prioriza la información actualizada y deja la réplica
como evolución posible. Esta es la única interacción síncrona del diseño;
su coste se analiza en la Sección~\ref{sec:p3-valoracion}.

\paragraph{Sin transacciones distribuidas.}
Emitir una póliza y crear su plan de cobro no forman una transacción atómica. Construir
transacciones entre servicios viola el principio de desacoplamiento, y la recomendación
es revisar la granularidad antes que coordinarlas; las sagas son la excepción, con
costes importantes \parencite[cap.~17]{richardsford2020}. El sistema acepta, por tanto,
la consistencia eventual: durante un breve intervalo, una póliza recién emitida puede no
tener todavía su plan de cobro.

\paragraph{Contratos publicados, no modelos.}
Los mensajes transportan identificadores y solo los datos que cada consumidor necesita, en
un formato publicado.
Nunca transportan objetos del modelo interno: cada consumidor traduce el mensaje a su
propio modelo (mecanismo M5 de la Sección~\ref{sec:p3-contaminacion}).

\subsection{Relación con una arquitectura orientada al dominio y el Adaptador de Servicio}
\label{sec:p3-arquitectura}

Los contextos delimitados organizan el sistema a gran escala. Internamente, cada contexto
se construye según la Arquitectura Cebolla (Onion) o la Arquitectura Limpia (Clean). El
diseño se apoya en tres principios. Los dos primeros son propios de estos estilos; el
tercero lo formula la Guía de Arquitectura N-Capas:
\begin{itemize}
  \item \textbf{El dominio en el centro.} El modelo de dominio ocupa el centro y todo el
        acoplamiento apunta hacia él. Las interfaces se definen en el núcleo y las
        implementaciones que dependen de tecnologías concretas viven en los bordes
        \parencite{palermo2008}.
  \item \textbf{La Regla de Dependencia.} Las dependencias del código fuente solo pueden
        apuntar hacia adentro, hacia las políticas de más alto nivel
        \parencite[203]{martin2017}.
  \item \textbf{El dominio desacoplado y una Capa de Aplicación delgada.} La Guía insiste en desacoplar la capa de Dominio
        mediante contratos, interfaces, inyección de dependencias e inversión de control,
        para que las reglas de negocio queden aisladas y soporten un ritmo alto de cambios
        \parencite[159]{delatorre2010}. La Capa de Aplicación debe ser delgada: coordina
        las tareas y las conversiones de formato, pero no contiene lógica de negocio
        \parencite[247]{delatorre2010}. Interpretamos que, en este pasaje, la Guía llama
        \enquote{lógica de negocio} a la lógica del dominio. En la terminología de Martin, los casos de uso sí
        contienen reglas de negocio propias de la aplicación, mientras que las reglas
        críticas residen en las entidades \parencite[204]{martin2017}.
\end{itemize}

La Figura~\ref{fig:p3-cebolla} aplica esta estructura al contexto de Siniestros.

\begin{figure}[!t]
  \centering
  \begin{tikzpicture}[
      capa/.style={draw=azul, rounded corners=6pt, thick},
      caja/.style={draw=azul, fill=white, rounded corners=2pt, font=\scriptsize\ttfamily,
                   inner sep=3pt, align=center},
      ext/.style={draw=gris, dashed, fill=gris!8, rounded corners=2pt, font=\scriptsize,
                  inner sep=4pt, align=center, minimum width=3.4cm},
      impl/.style={-{Triangle[open,length=2.4mm,width=2.2mm]}, thick, azul},
      usa/.style={-{Stealth[length=2.2mm]}, thick, gris},
      tit/.style={font=\scriptsize\bfseries, azul}
    ]
    % Capas concéntricas (de afuera hacia adentro)
    \draw[capa, fill=azul!4]  (-7.8,-4.9) rectangle (7.8,4.1);
    \draw[capa, fill=azul!10] (-7.45,-2.65) rectangle (7.45,2.35);
    \draw[capa, fill=azul!18] (-5.4,-2.15) rectangle (5.4,1.75);
    \draw[capa, fill=azul!30] (-3.4,-0.55) rectangle (3.4,1.05);

    \node[tit] at (0,3.75) {Infraestructura: adaptadores (detalles tecnológicos)};
    \node[tit] at (0,2.05) {Aplicación (casos de uso): \mdseries AbrirSiniestro · SolicitarPeritaje · RegistrarPeritaje · LiquidarSiniestro};
    \node[tit] at (0,1.40) {Servicios de dominio y puertos (interfaces del núcleo)};
    \node[font=\small, align=center] at (0,0.25)
      {\textbf{Modelo de dominio}\\[1pt]\footnotesize Siniestro (raíz del agregado) · CoberturaAplicable};

    % Puertos definidos en el núcleo
    \node[caja] (p1) at (-3.45,-1.35) {IConsultaCobertura};
    \node[caja] (p2) at (0,-1.35)     {IRepositorioSiniestros};
    \node[caja] (p3) at (3.45,-1.35)  {IPublicadorEventos};

    % Adaptadores de salida (implementan los puertos)
    \node[caja, fill=naranja!10, draw=naranja] (a1) at (-4.7,-3.65)
      {AgenteServicioSuscripcion\\\normalfont(Agente de Servicio)};
    \node[caja] (a2) at (0,-3.65)   {RepositorioSiniestrosSQL};
    \node[caja] (a3) at (4.7,-3.65) {PublicadorBroker};

    % Adaptadores de entrada
    \node[caja] (e1) at (-3.6,3.05) {ConsumidorEventos (PeritajeCompletado)};
    \node[caja] (e2) at (3.6,3.05)  {UI de tramitadores};

    % Sistemas externos
    \node[ext] (x1) at (-4.7,-6.0) {API publicada de\\Suscripción y Pólizas};
    \node[ext] (x2) at (0,-6.0)    {Base de datos propia\\de Siniestros};
    \node[ext] (x3) at (4.7,-6.0)  {Bróker de mensajes};

    \draw[impl] (a1) -- (p1);
    \draw[impl] (a2) -- (p2);
    \draw[impl] (a3) -- (p3);
    \draw[usa]  (a1) -- (x1);
    \draw[usa]  (a2) -- (x2);
    \draw[usa]  (a3) -- (x3);
    \draw[usa]  (e1.south) -- ++(0,-0.62);
    \draw[usa]  (e2.south) -- ++(0,-0.62);

    % Leyenda
    \draw[impl] (-7.6,-7.15) -- ++(0.9,0) node[right,font=\scriptsize,black]
      {implementa un puerto (la dependencia apunta hacia el núcleo)};
    \draw[usa]  (3.4,-7.15) -- ++(0.9,0) node[right,font=\scriptsize,black]
      {usa / invoca};
  \end{tikzpicture}
  \caption{Estructura interna Onion/Clean del contexto de Siniestros. El Agente de
           Servicio es el único componente que conoce el protocolo y el formato de
           Suscripción.}
  \label{fig:p3-cebolla}
\end{figure}

\paragraph{El Adaptador de Servicio.}
Lo que el enunciado llama \enquote{Adaptador de Servicio} corresponde en la Guía a los
\emph{Agentes de Servicios externos}.\footnote{La Guía emplea además
\enquote{Agente de Servicio} en otro sentido: un componente de la capa de presentación
que consume servicios distribuidos \parencite[269]{delatorre2010}. En esta respuesta,
\enquote{Agente de Servicio} designa siempre a un agente de servicios externos.} Según la Guía, estos agentes:
\begin{itemize}
  \item gestionan la semántica de comunicación con un servicio externo y mapean entre su
        formato de datos y el de la aplicación \parencite[108]{delatorre2010};
  \item aíslan las particularidades de ese servicio, de modo que, manteniendo ciertas
        interfaces, el servicio puede sustituirse por otro sin afectar al sistema
        \parencite[58]{delatorre2010};
  \item en aplicaciones grandes actúan como nivel de abstracción entre la capa de
        Dominio y los servicios remotos, con una interfaz homogénea. Conviene
        desacoplarlos mediante inversión de control, lo que permite simular los servicios
        externos con \textit{fakes} durante el desarrollo y probar los propios agentes de
        forma unitaria \parencite[155--156]{delatorre2010}.
\end{itemize}

Es la misma idea que el adaptador de la Arquitectura Hexagonal
\parencite{cockburn2005} y que la capa de adaptadores de interfaz de la Arquitectura
Limpia, que incluye los adaptadores hacia servicios externos \parencite[205]{martin2017}.
En el diseño propuesto, \texttt{AgenteServicioSuscripcion} implementa el puerto
\texttt{IConsultaCobertura}, definido en el núcleo de Siniestros. Llama a la API
publicada de Suscripción y traduce la respuesta al objeto de valor local
\texttt{CoberturaAplicable} (Listado~\ref{lst:p3-agente}).

\begin{lstlisting}[style=codigo, language={[Sharp]C}, float=!tb,
  caption={Esbozo en estilo C\# del puerto y del Agente de Servicio en el contexto de Siniestros.},
  label={lst:p3-agente}]
// --- Núcleo de Siniestros: el puerto lo define el dominio ---
public interface IConsultaCobertura {
    CoberturaAplicable ObtenerCobertura(string numeroPoliza, DateTime fechaSiniestro);
}

// Objeto de valor propio de Siniestros: NO es la Póliza de Suscripción
public record CoberturaAplicable(string NumeroPoliza, bool Vigente,
                                 decimal LimiteIndemnizacion, decimal Franquicia);

// --- Infraestructura de Siniestros: Agente de Servicio (adaptador) ---
public class AgenteServicioSuscripcion : IConsultaCobertura {
    private readonly IClienteHttp http;   // el detalle tecnológico queda confinado aquí

    public AgenteServicioSuscripcion(IClienteHttp http) { this.http = http; }

    public CoberturaAplicable ObtenerCobertura(string numeroPoliza, DateTime fecha) {
        // Contrato publicado por Suscripción (formato externo)
        string ruta = $"/polizas/{numeroPoliza}/cobertura?fecha={fecha:yyyy-MM-dd}";
        CoberturaDto dto = http.Get<CoberturaDto>(ruta);
        // Traducción al modelo de Siniestros: lo externo no entra al núcleo
        return new CoberturaAplicable(dto.Numero, dto.Estado == "VIGENTE",
                                      dto.LimiteDanos, dto.Franquicia);
    }
}
\end{lstlisting}

El Agente de Servicio cierra el círculo de la Sección~\ref{sec:p3-contaminacion}. Es la
barrera física que impide que el modelo de Póliza de Suscripción entre en el núcleo de
Siniestros, y aporta tres beneficios:
\begin{itemize}
  \item un cambio en la API o en el formato de Suscripción solo afecta a este adaptador;
  \item el dominio de Siniestros puede probarse de forma aislada con una implementación
        simulada del puerto, del mismo modo que la Arquitectura Hexagonal permite
        desarrollar y probar la aplicación aislada de sus dispositivos y bases de datos
        \parencite{cockburn2005};
  \item como el núcleo solo conoce la interfaz \texttt{IConsultaCobertura}, se respeta la
        Regla de Dependencia \parencite[203]{martin2017}.
\end{itemize}
Los consumidores de eventos cumplen el mismo papel de traducción en sentido de entrada:
convierten cada mensaje publicado en una llamada a un caso de uso del contexto receptor.

\subsection{Valoración crítica}
\label{sec:p3-valoracion}

La solución reúne las ventajas que justifican DDD en un dominio con reglas de negocio
cambiantes y varias áreas de expertos:
\begin{itemize}
  \item cada modelo evoluciona al ritmo de su área;
  \item los cambios quedan contenidos dentro de un contexto;
  \item el lenguaje del código coincide con el de los expertos;
  \item cada núcleo de dominio puede probarse sin infraestructura.
\end{itemize}

Tiene, sin embargo, costes que deben asumirse de forma consciente:
\begin{enumerate}[label=(\roman*)]
  \item \textbf{Duplicación de datos y de código entre contextos}, que es el precio del
        desacoplamiento \parencite[cap.~17]{richardsford2020}.
  \item \textbf{Consistencia eventual y menor control del flujo.} Los contextos pueden
        estar momentáneamente desalineados. Además, la topología \textit{broker} complica
        el control del flujo, la gestión de errores y la recuperación, porque ningún
        componente supervisa el proceso completo \parencite[cap.~14]{richardsford2020}.
  \item \textbf{Dependencia temporal en la consulta síncrona.} Si Suscripción no está
        disponible, Siniestros no puede verificar la cobertura al abrir un expediente. Es
        el mismo riesgo por el que se evitó acoplar de forma síncrona a Peritaje
        (criterio C4). Como mitigación, propia de este diseño, Siniestros puede registrar
        el expediente en estado \enquote{cobertura pendiente de verificar} y completar la
        verificación cuando Suscripción responda, sin bloquear la atención al asegurado. Para
        ello, el puerto \texttt{IConsultaCobertura} debe contemplar el fallo de la
        llamada; el Listado~\ref{lst:p3-agente} omite ese manejo por brevedad.
  \item \textbf{Más infraestructura y riesgo de granularidad excesiva.} Hay que operar un
        bróker de mensajes. Además, el error frecuente de crear servicios demasiado
        pequeños obliga después a reconstruir enlaces de comunicación entre ellos
        \parencite[cap.~17]{richardsford2020}.
\end{enumerate}

La propia Guía advierte que, si una aplicación es sencilla, sus reglas de negocio
cambiarán poco y no se prevén cambios de infraestructura, una arquitectura orientada al
dominio probablemente no compensa \parencite[42]{delatorre2010}. No es el caso de la
aseguradora: cuatro áreas con lenguajes propios y reglas contractuales, de cobro y de
peritaje que evolucionan de forma independiente son justamente el escenario en el que la
división en contextos delimitados, comunicados por eventos y protegidos por Agentes de
Servicios externos, aporta más valor que coste.
